\documentclass[11pt,a4paper]{article}
\usepackage{astra-paper}

\newcommand{\CW}{\mathrm{CW}}
\newcommand{\Shapes}{\mathcal S}
\DeclareMathOperator{\Split}{Split}
\DeclareMathOperator{\supp}{supp}
\DeclareMathOperator{\rk}{rk}
\DeclareMathOperator{\PH}{PH}
\DeclareMathOperator{\KL}{KL}
\newcommand{\brk}{\underline{\rk}}
\newcommand{\dg}{\mathrel{\ge_{\mathrm{deg}}}}
\newcommand{\Tpos}[1]{\mathcal T_{#1}^{+}}
\newcommand{\Etot}{E_{\mathrm{total}}}
\newcommand{\Mtot}{M_{\mathrm{total}}}
\newcommand{\Eshared}{E^{\mathrm{shared}}}
\newcommand{\Mbar}{\overline M}
\newcommand{\Ngood}{N_{\mathrm{good}}}
\newcommand{\Nadm}{N_{\mathrm{adm}}}
\newcommand{\RR}{\mathbb R}
\newcommand{\ZZ}{\mathbb Z}
\newcommand{\WW}{\{X,Y,Z\}}
\newcommand{\abs}[1]{\left|#1\right|}
\newcommand{\norm}[1]{\left\lVert#1\right\rVert}
\newcommand{\inner}[2]{\left\langle#1,#2\right\rangle}
\newcommand{\eps}{\varepsilon}

\begin{document}
\paperTitle

\begin{abstract}
\noindent We improve the matrix multiplication exponent to $\omega<2.371045$, compared with the bound $\omega<2.371177$ of Dupont et al.\ (2026). Our main contribution is a change to the extraction step of the asymmetric laser method. Earlier analyses extract each recursion level separately, passing up cross-level optimizations. We prove that factors from different levels can share one asymmetric hash. A finite, staggered pipeline makes these factors available simultaneously and combines their contributions before taking the minimum over axes, up to an explicit boundary loss. We formulate the resulting optimization problem in the notation of Dupont et al.\ and,  using automatic differentiation and constrained quasi-Newton optimization, converge to this bound. The optimized parameters are rounded to exact rationals and certified independently of the search. The entire argument, from the Coppersmith--Winograd border rank through the shared extraction, the pipeline and the asymptotic sum inequality to the numerical certificate, is formalized in Lean~4 as a statement about exact tensor rank; the formal development has been rebuilt independently and its statement audited. The construction and its proof were produced largely by AI systems under our direction, and we rely on the formal proof for correctness.
\end{abstract}
\vspace{2pt}
\noindent\textbf{Keywords.}\ Matrix multiplication; laser method; asymmetric hashing; numerical optimization; formal verification.

\section{Preface}\label{sec:preface}

We are not laser-method specialists. Our investigation began with two questions: can current AI systems improve on the bound of \citet{dupont2026}, and can a laser-method proof be formalized end to end? Community efforts such as Prove2Me \citep{prove2me2026} are formalizing the classical bounds on $\omega$ mission by mission against human-confirmed definitions. Beginning on 17 September 2026 we prompted OpenAI's GPT 6 Astra Pro via simple prompting to "make significant analytical progress on [improving the bound]". After about 4 hours of work it produced a natural language proof. Similarly simple prompting produced the kernel-checked theorem for the final bound after around 30 hours of work. See Appendix~\ref{sec:timeline} for the full timeline. The formalization used two consumer subscriptions (three usage resets of an OpenAI Codex Pro plan and about half of an Anthropic Claude Max allowance) and comprises $88{,}134$ declarations. Although we certify the theorem statement (120 Lean loc) and the Lean verification, we cannot speak to every step of the proof. Turning machine-checked proofs of this scale into clear exposition is, in our estimation, an open and existential problem moving forward in the fields of mathematics and computer science research. We hope that experts on the laser method will find the shared-extraction idea worthy of further investigation. In conjunction with this paper, we release a \href{https://oliverproudfoot.substack.com/p/pocket-sized-super-mathematical-intelligence}{blog post} with additional thoughts. The following writeup was generated by GPT-6 Astra (Pro) with stylistic feedback from the authors.

\section{Introduction}\label{sec:Intro}

The best known constructions for fast matrix multiplication combine two kinds of progress: a mathematical analysis that turns tensor decompositions into algorithms, and a numerical search for parameters that make that analysis effective. Improving the search can strengthen an existing bound. Changing the analysis can make a different set of parameters useful.

The Coppersmith--Winograd (\(\CW\)) tensor is central to this approach \citep{cw1990}. Recent advances use asymmetric hashing to reduce the combination loss incurred when extracting matrix multiplication tensors from its powers \citep{dwz2023,vxxz2024}. \citet{alman2025} further develop this analysis by treating all three axes asymmetrically. \citet{dupont2026} give an explicit optimization formulation of that analysis and use modern optimization and AlphaEvolve to obtain $\omega<2.371177$.

The significance of $\omega$ extends beyond the multiplication primitive itself. Classical reductions transfer improvements in matrix multiplication to dense inversion, linear-system solution, and least-squares computation \citep{demmel2007}; more broadly, $\omega$ is the asymptotic-rank exponent of the matrix multiplication tensor and a benchmark for the reach of current tensor-decomposition methods.

Our contribution is to change \emph{which factors are extracted together}. In a level-by-level analysis, the number of copies retained at each level is constrained by that level's least favorable axis. These limiting axes can differ across levels. A joint extraction can exploit their complementarity when the factors are present at the same time and the compatibility analysis applies to their product.

Our heterogeneous extraction theorem establishes the required joint step. Its hash remains valid when parent positions have different word lengths, and its compatibility counts add across the different factor types. A staggered pipeline then applies the theorem to factors from levels $4$, $3$ and $2$ belonging to different source batches. Every factor descends exactly once per round, and all start-up and drain costs are included.

\subsection{Our result and contribution}

For a commutative ring $R$, let $\rk_R(\langle n,n,n\rangle)$ denote the minimum number of terms in an exact rank decomposition of the matrix multiplication tensor, and define
\begin{equation}\label{eq:omega}
\omega_R=\inf\!\left\{\tau\ge0:\ \exists c>0\ \forall n\ge1,\quad
\rk_R(\langle n,n,n\rangle)\le c n^\tau\right\}.
\end{equation}
Over fields, the usual connection between tensor rank and bilinear algorithms identifies this exponent with the arithmetic matrix multiplication exponent \citep{bcs1997,blaser2013}. We write $\omega$ when the ground field is understood.

\begin{resultbox}
\begin{theorem}[Main result]\label{thm:main}
For every commutative ring $R$ in the universe of the accompanying formal statement,
\[
\omega_R<2.371045.
\]
In particular, over fields this yields the corresponding bound on the usual matrix multiplication exponent.
\end{theorem}
\end{resultbox}

The mathematical exposition below uses a characteristic-zero field. The formal development proves the ring-theoretic exact-rank statement in \cref{eq:omega}; its scope is recorded in \cref{app:formal}.

\paragraph{A new extraction theorem.}
We extend one asymmetric extraction to a fixed finite product of interface factors at different recursion levels. The local entropy formulas are those of \citet{alman2025}. The new theorem allows their contributions to be combined \emph{before} taking the minimum over the three hashing roles.

\paragraph{A finite realization.}
We use six coordinate-permuted sources to align the hashing roles, and a $K$-batch pipeline to respect dependencies between levels. The final retained rate includes the start-up and drain costs through an explicit correction $D/K$.

\paragraph{An optimized, certifiable witness.}
The changed analysis leads to a different nonconvex optimization problem. A constrained, automatic-differentiation-based search finds parameters for $\CW_5^{\otimes8}$. Rationalization and a separate certificate establish the numerical inequality. The accompanying Lean development supplies an additional end-to-end check of the exact-rank theorem.

\begin{table}[t]
\centering\small
\begin{tabularx}{\linewidth}{@{}l l X@{}}
\toprule
Work & Bound & Principal change\\
\midrule
\citet{dwz2023} & $2.371866$ & Asymmetric hashing\\
\citet{vxxz2024} & $2.371552$ & Refined combination-loss analysis\\
\citet{alman2025} & $2.371339$ & Additional asymmetry between axes\\
\citet{dupont2026} & $2.371177$ & Larger-scale parameter optimization\\
\textbf{This work} & $\mathbf{2.371045}$ & \textbf{Shared extraction across recursion levels}\\
\bottomrule
\end{tabularx}
\caption{Recent upper bounds. The value shown for \citet{alman2025} is the original SODA~2025 bound. Our contribution combines a shared extraction theorem with optimization of the resulting formulation.}
\label{tab:history}
\end{table}

\paragraph{Concurrent and related work.}
\citet{naslund2026} publicly shared a transcript in which the same AI system, prompted to improve the bound of \citet{dupont2026}, reports $\omega<2.371161$ with a standalone interval-arithmetic verifier. That witness stays within the level-by-level formulation \cref{eq:old-opt}: it is a better point of the same optimization problem, found with optimization techniques close to those of \cref{sec:optimization}. The present bound instead changes the extraction theorem, and in \cref{app:search} we report how far our own searches reached without that change. On the formal side, the Prove2Me platform \citep{prove2me2026} hosts a community campaign to formalize upper bounds on $\omega$ in Lean against human-confirmed definitions. Our development is independent of that effort; relating our rank-based definition \cref{eq:omega} to the operation-count definition used there is natural future work.

\section{The optimization problem and its notation}\label{sec:model}

We retain the notation of \citet[Section~2]{dupont2026}: the root $G$, nodes $T$, shapes $s_T$, regions $r$ with permutations $\pi_r$, masses $m_T$, splits $\alpha_T^{(r)}$, complete split distributions $\beta_{T,W}$, retained exponents $E$, and logarithmic matrix sizes $M$. The new quantities will be the role sums $e_{\ell,W}$, the shared rate $\Eshared$, and the pipeline correction $D/K$.

Throughout, logarithms and entropies are base two. The integers $q\ge1$ and $\ell^*\ge2$ are fixed hyperparameters. Our witness uses
\begin{equation}\label{eq:source}
q=5,\qquad \ell^*=4,\qquad \CW_5^{\otimes8},\qquad
C=2^{\ell^*-1}\log(q+2)=8\log7.
\end{equation}
The cost $C$ is the logarithmic border-rank budget per source copy. With the convention that a level-$\ell$ shape has coordinate sum $2^\ell$, a level-$\ell$ constituent contains $2^{\ell-1}$ base CW factors.

\subsection{Tree structure and free parameters}\label{sec:tree}

A level-$\ell$ shape is
\[
\Shapes_\ell=\{(i,j,k)\in\ZZ_{\ge0}^3:i+j+k=2^\ell\}.
\]
A non-root node is \emph{positive} if all three coordinates of $s_T$ are positive; otherwise it is a \emph{zero-shape} node. There are six regions, with $\pi_r$ the $r$th permutation of $X,Y,Z$ in lexicographic order. Region $r$ processes the axes in the order $\pi_r(X),\pi_r(Y),\pi_r(Z)$.

For each $s\in\Shapes_{\ell^*}$ and $r\in[6]$, the root has a child $G[s,r]$. A positive node of level $\ell\ge3$ has children $T[u,r]$, where
\begin{equation}\label{eq:split}
\Split(s_T)=\{u\in\Shapes_{\ell-1}:u\le s_T\}
\end{equation}
contains the possible left-child shapes. Its right-child shape is $s_T-u$. Zero-shape nodes and positive level-$2$ nodes are leaves; their set is $\mathcal L$. We write $\Tpos\ell$ for the positive nodes at level $\ell$.

For a finite set $D$, $\Delta(D)$ denotes its probability simplex. The root and each positive node of level at least three carry region weights $A_T\in\Delta([6])$ and split distributions $\alpha_T^{(r)}$. The root splits lie in $\Delta(\Shapes_{\ell^*})$; all other splits lie in $\Delta(\Split(s_T))$.

A complete split distribution records the full words defining fine variable blocks. Define
\[
\mathcal C_{\ell,a}=\left\{L\in\{0,1,2\}^{2^{\ell-1}}:\sum_pL_p=a\right\}.
\]
A distribution on $\mathcal C_{\ell,a}$ specifies the frequency of each such word. At a zero-shape node, the distribution on the first nonzero coordinate is free. The zero coordinate has its all-zero word, and the third distribution is determined by the complement map $L\mapsto\vec2-L$.

A positive level-$2$ node has shape $(1,1,2)$ or a permutation thereof and a parameter $\mu_T\in[0,1/2]$. For shape $(1,1,2)$,
\begin{equation}\label{eq:leveltwo-beta}
\begin{aligned}
\beta_{T,X}=\beta_{T,Y}&=\tfrac12\delta_{0,1}+\tfrac12\delta_{1,0},\\
\beta_{T,Z}&=\mu_T\delta_{0,2}+\mu_T\delta_{2,0}+(1-2\mu_T)\delta_{1,1}.
\end{aligned}
\end{equation}
Here $\delta_{a,b}$ is the point mass at the two-letter word $(a,b)$.

\subsection{Masses and complete split distributions}

Masses count copies per source copy. They are propagated from the root downward:
\begin{gather}
m_G=1,\qquad m_{G[s,r]}=A_G^{(r)}\alpha_G^{(r)}(s),\label{eq:root-mass}\\
m_{T[u,r]}=m_TA_T^{(r)}
\underbrace{\bigl(\alpha_T^{(r)}(u)+\alpha_T^{(r)}(s_T-u)\bigr)}_{\text{\upshape both halves}}.\label{eq:mass}
\end{gather}
The two terms in \cref{eq:mass} account for a child occurring in either half of a parent.

Complete distributions are propagated in the opposite direction. For a positive node,
\begin{align}
\beta_{T,W}^{(r)}
&=\sum_{u\in\Split(s_T)}\alpha_T^{(r)}(u)
\bigl(\beta_{T[u,r],W}\times\beta_{T[s_T-u,r],W}\bigr),\label{eq:beta-region}\\
\beta_{T,W}&=\sum_{r=1}^6 A_T^{(r)}\beta_{T,W}^{(r)}.\label{eq:beta}
\end{align}
The product $\times$ concatenates two independently drawn words. Thus the numerical evaluator has a bottom-up pass for word distributions and a top-down pass for multiplicities. The free parameters determine both passes completely.

\subsection{The local retained exponents}

For $\rho\in\Delta(D)$, let $\rho_W$ be its marginal in coordinate $W$, and define
\begin{align}
H_D^{\max}(\rho)&=\sup\{H(\rho'): \rho'\in\Delta(D),\ \rho'_W=\rho_W\text{ for all }W\},\label{eq:hmax}\\
P_D(\rho)&=H_D^{\max}(\rho)-H(\rho).\label{eq:penalty}
\end{align}
The penalty measures the multiplicity of other joint split types with the same marginals.

For a positive node $T$ of level at least three, the retained exponents in the identity region are
\begin{equation}\label{eq:local-E}
\begin{aligned}
E_{T,X}^{(1)}&=m_TA_T^{(1)}
\left[H\bigl((\alpha_T^{(1)})_X\bigr)-P_{\Split(s_T)}(\alpha_T^{(1)})\right],\\
E_{T,Y}^{(1)}&=m_TA_T^{(1)}\left[H(\beta_{T,Y}^{(1)})-\eta_{T,Y}^{(1)}\right],\\
E_{T,Z}^{(1)}&=m_TA_T^{(1)}\left[H(\beta_{T,Z}^{(1)})-\eta_{T,Z}^{(1)}\right].
\end{aligned}
\end{equation}
The quantities $\eta_{T,Y}^{(1)}$ and $\eta_{T,Z}^{(1)}$ are the pooled compatibility entropies: they bound how many fine blocks remain compatible with a competing coarse triple. Their explicit formulas are given in \cref{app:local}. Other regions are obtained by relabelling the axes according to $\pi_r$.

In the level-by-level formulation,
\begin{equation}\label{eq:old-level}
E_\ell=\sum_{r=1}^6\min_{W\in\WW}
\sum_{T\in\Tpos\ell}E_{T,W}^{(r)},\qquad \ell\ge3.
\end{equation}
Each region of each level therefore pays for its own least favorable axis. At level two, the local vector for shape $(1,1,2)$ is
\[
(E_{T,X},E_{T,Y},E_{T,Z})
=m_T\bigl(1,1,H(\mu_T,\mu_T,1-2\mu_T)\bigr),
\]
with coordinates permuted for the other shapes; $E_2=\min_W\sum_{T\in\Tpos2}E_{T,W}$.

The root is extracted separately. Its rate $E_G$ is the standard root rate of the same formulation and enters the final bound as a separate contribution. Its local entropy formulas are recorded in \cref{app:local}.

\subsection{The original final constraint}

Each leaf supplies a rectangular matrix multiplication tensor with logarithmic side contributions $M_{T,X},M_{T,Y},M_{T,Z}$. These expressions are unchanged in our construction and are stated in \cref{app:local}. The level-by-level analysis of \citet{alman2025} assembles the quantities as
\begin{equation}\label{eq:old-total}
\Etot=E_G+E_2+\sum_{\ell=3}^{\ell^*}E_\ell,
\qquad
\Mtot=\min_{W\in\WW}\sum_{T\in\mathcal L}M_{T,W}.
\end{equation}
The optimization problem is
\begin{equation}\label{eq:old-opt}
\begin{aligned}
\operatorname{minimize}\quad &\Omega\\[3pt]
\text{\normalfont subject to}\quad &\Etot+\Omega\Mtot\ge2^{\ell^*-1}\log(q+2).
\end{aligned}
\end{equation}
Any feasible solution yields an upper bound on $\omega$ \citep{alman2025,dupont2026}. We retain the tree, local distributions, entropy penalties and leaf dimensions. The change is in the assembly of the constituent retention rates.

\section{Shared extraction across recursion levels}\label{sec:shared}

\subsection{Loss from level-separated extraction}

It is useful to distinguish a physical axis from a hashing role. In region $r$, role $W$ is played by physical axis $\pi_r(W)$. Define the role-summed level contributions
\begin{equation}\label{eq:role-sums}
e_{\ell,W}=\sum_{r=1}^6\sum_{T\in\Tpos\ell}
E_{T,\pi_r(W)}^{(r)},\qquad \ell=3,4.
\end{equation}
For a terminal node, let $v_T\in\Delta(\WW)$ specify which hashing role receives its doubled coordinate. Writing $H_T=H(\mu_T,\mu_T,1-2\mu_T)$, its contribution is
\begin{equation}\label{eq:terminal-role}
e_{2,W}=\sum_{T\in\Tpos2}m_T\bigl[1+(H_T-1)v_{T,W}\bigr].
\end{equation}
These assignments are realized by partitioning positions into sectors, as described in \cref{app:roles}.

For a fixed role allocation, level-separated extraction pays
$\sum_{\ell=2}^4\min_W e_{\ell,W}$. Our shared rate is
\begin{resultbox}
\begin{equation}\label{eq:shared}
\Eshared=\min_{W\in\WW}
\underbrace{\bigl(e_{4,W}+e_{3,W}+e_{2,W}\bigr)}_{\text{\upshape pooled levels}}.
\end{equation}
\end{resultbox}
The inequality $\Eshared\ge\sum_\ell\min_W e_{\ell,W}$ identifies the potential gain. The extraction theorem below constructs a tensor degeneration attaining the shared rate in \cref{eq:shared}.

\begin{figure}[t]
\centering

\begin{minipage}[t]{.478\linewidth}
\centering
\vspace*{-2.5pt}
{\bfseries\small Level-separated extraction}\par\vspace{7pt}
{\footnotesize\setlength{\tabcolsep}{4pt}\renewcommand{\arraystretch}{1.35}
\begin{tabular}{@{}lccc@{}}
\toprule
 & $X$ & $Y$ & $Z$\\
\midrule
Level 4 & 1.6637 & 1.6451 & \cellcolor{wash}\textbf{1.6313}\\
Level 3 & 2.2678 & 2.2081 & \cellcolor{wash}\textbf{1.8906}\\
Level 2 & \cellcolor{wash}\textbf{1.5094} & 1.5877 & 1.9190\\
\bottomrule
\end{tabular}}\par\vspace{8pt}
{\small Select each row's minimum, then add.}\par\vspace{5pt}
{\color{muted}$\displaystyle\sum_\ell\min_W e_{\ell,W}=$}\par\vspace{3pt}
{\fontsize{17}{20}\selectfont\bfseries\color{muted}5.031390}
\end{minipage}\hfill
\begin{minipage}[t]{.478\linewidth}
\centering
{\bfseries\small\color{accent} Shared extraction}\par\vspace{7pt}
{\footnotesize\setlength{\tabcolsep}{4pt}\renewcommand{\arraystretch}{1.35}
\begin{tabular}{@{}lccc@{}}
\toprule
 & $X$ & $Y$ & $Z$\\
\midrule
Level 4 & 1.6637 & 1.6451 & 1.6313\\
Level 3 & 2.2678 & 2.2081 & 1.8906\\
Level 2 & 1.5094 & 1.5877 & 1.9190\\
\bottomrule
\end{tabular}}\par\vspace{8pt}
{\small Add each column, then take the minimum.}\par\vspace{5pt}
{\color{accent}$\displaystyle\min_W\sum_\ell e_{\ell,W}=$}\par\vspace{3pt}
{\fontsize{17}{20}\selectfont\bfseries\color{accent}5.440956}
\end{minipage}

\caption{\textbf{The numerical opportunity.} At the supplied witness, level~3 is least favorable in role $Z$, whereas level~2 is least favorable in role $X$. Taking row minima loses their complementarity. A shared extraction first adds the columns, which the search nearly equalizes. Both sides evaluate the displayed witness with the same parameters and role allocations.}
\label{fig:balance}
\end{figure}

The distinction is visible in \cref{fig:balance}. The sum of the separate level minima is approximately $5.031390$ bits per source copy. The shared rate is $5.440956$ bits. Pooling the three levels recovers approximately $0.409566$ bits at this witness. The original region-by-region rates can be smaller still: \cref{eq:old-level} takes its minima before summing over regions.

\subsection{A heterogeneous extraction theorem}

An \emph{interface} is a product of constituent tensors restricted to prescribed empirical distributions of their fine words. Its tolerance and complete-word constraints specify the input required by the next extraction.

\begin{theorem}[Shared extraction; informal form]\label{thm:hetero-informal}
Fix finitely many interface-factor types $t$, possibly at different recursion levels. Let $m_t$ be their multiplicities per source copy, and let $e_{t,X},e_{t,Y},e_{t,Z}$ be their usual local retained exponents per unit mass, in one common role order. A joint asymmetric extraction retains
\[
2^{N(g-\rho)-o(N)}\quad\text{copies},\qquad
 g=\min_W\sum_t m_t e_{t,W},
\]
of the prescribed child product, for every fixed $\rho>0$, with appropriately chosen tolerances and subexponential overhead.
\end{theorem}

The precise quantifiers and interface definitions appear in \cref{app:theorem}. Three points require proof. Different word lengths must share the same hash; compatibility tests must remain valid variable restrictions; and the surviving child products must be repaired and thickened so that the next extraction can use them.

\paragraph{A shared hash for mixed levels.}
At parent position $p$, let $(I_p,J_p,K_p)$ be the left-child shape. If that parent has level $\ell_p$, then
\[
I_p+J_p+K_p=c_p,\qquad c_p=2^{\ell_p-1}.
\]
Choose an odd prime modulus $M$ and independent uniform coefficients $b_0,w_0,w_p$ modulo $M$. Define
\begin{equation}\label{eq:hash}
\begin{aligned}
h_X(I)&=b_0+\sum_p w_p I_p,\\
h_Y(J)&=b_0+w_0+\sum_p w_p J_p,\\
h_Z(K)&=b_0+\tfrac12\left(w_0+\sum_p w_p(c_p-K_p)\right).
\end{aligned}
\end{equation}
Then $h_X+h_Y=2h_Z$ on every admissible triple, even though $c_p$ varies with the level. Restricting the hashes to a progression-free set forces each surviving triple into a common bucket. Conditional on one triple occupying a specified bucket, a distinct triple sharing one of its coarse vertices collides with probability $1/M$.

\paragraph{One modulus controls all factor types.}
Let $\Nadm$ count admissible coarse triples, $\Ngood$ the triples with the prescribed split types, $B_W$ the coarse vertices on axis $W$, and $p_Y,p_Z$ the fractions compatible with a fixed typical fine block. The single modulus is chosen against
\begin{equation}\label{eq:modulus}
M_0=\left\lceil N^8\max\left\{1,\frac{\Nadm}{B_X},
\frac{\Ngood p_Y}{B_Y},\frac{\Ngood p_Z}{B_Z}\right\}\right\rceil,
\qquad M\in[M_0,2M_0].
\end{equation}
The counts factor over the fixed type-position sets. Their logarithms therefore add. Subtracting $\log M$ from $\log\Ngood$ gives
\begin{equation}\label{eq:rate-mechanism}
\log\Ngood-\log M
\ge N\min_W\sum_t m_t e_{t,W}-Nr(\eps)-o(N),
\end{equation}
where $r(\eps)\to0$. This is the point at which the minimum moves outside the type sum.

The compatibility zero-outs are performed in the order $X$, then $Y$, then $Z$. Parent-type concentration leaves only a small fraction of holes in the intended children. A partitioned-tensor repair argument costs $2^{o(N)}$ copies, and a polynomial-size decomposition into exact empirical types produces the approximate interfaces needed next. \Cref{app:proof} proves the complete sequence of restrictions, repair, and interface conversion.

\subsection{Making the levels available at the same time}\label{sec:pipeline}

Each child becomes available after its parent has been extracted. Staggering \emph{different batches} brings factors from several levels into the same round. First use six coordinate-permuted source copies and regroup strategies by physical hashing order. Every physical region then has the same role-summed vector; applying \cref{thm:hetero-informal} once in each region yields \cref{eq:shared} after normalization. The root rate remains a separate contribution.

Fix $K\ge2$ before $N$ grows. After separate root extractions, round $j$ processes batch $j$ at level~4, batch $j-1$ at level~3, and batch $j-2$ at level~2, omitting nonexistent batches. The steady rounds share a hash across all three active levels. The first two rounds fill the pipeline; the last two drain it.

\begin{figure}[t]
\centering

{\small\setlength{\tabcolsep}{6.3pt}\renewcommand{\arraystretch}{1.35}
\begin{tabular}{@{}l>{\columncolor{wash}}c>{\columncolor{wash}}cccc>{\columncolor{wash}}c>{\columncolor{wash}}c@{}}
\toprule
Round & $1$ & $2$ & $3$ & $\cdots$ & $K$ & $K+1$ & $K+2$\\
\midrule
Level 4 & $1$ & $2$ & $3$ & $\cdots$ & $K$ & -- & --\\
Level 3 & -- & $1$ & $2$ & $\cdots$ & $K-1$ & $K$ & --\\
Level 2 & -- & -- & $1$ & $\cdots$ & $K-2$ & $K-1$ & $K$\\
\midrule
Retained rate & $g_4$ & $g_{43}$ & $g_{432}$ & $\cdots$ & $g_{432}$ & $g_{32}$ & $g_2$\\
\bottomrule
\end{tabular}}
\par\vspace{6pt}

\caption{\textbf{The finite pipeline.} Entries are batch indices. Each batch appears exactly once at each level. The shaded boundary rounds are retained in the analysis; all other rounds achieve $g_{432}=\Eshared$. The number $K$ is fixed before the tensor-power parameter $N$ tends to infinity.}
\label{fig:pipeline}
\end{figure}

Write
\begin{equation}\label{eq:round-rates}
\begin{gathered}
g_4=\min e_4,\qquad g_{43}=\min(e_4+e_3),\qquad
 g_{432}=\min(e_4+e_3+e_2)=\Eshared,\\
g_{32}=\min(e_3+e_2),\qquad g_2=\min e_2,
\end{gathered}
\end{equation}
where each minimum is over three coordinates. Reading \cref{fig:pipeline} by columns gives the exact total
\begin{align}
\underbrace{g_4+g_{43}}_{\text{\upshape start-up}}
+\underbrace{(K-2)g_{432}}_{\text{\upshape steady}}
+\underbrace{g_{32}+g_2}_{\text{\upshape drain}}
&=K\Eshared-D,\label{eq:pipeline-total}\\
D&=2g_{432}-g_4-g_{43}-g_{32}-g_2\ge0.\label{eq:boundary}
\end{align}
Indeed, $g_4+g_{32}\le g_{432}$ and $g_{43}+g_2\le g_{432}$. The loss per batch is exactly $D/K$.

Inactive factors are carried unchanged between rounds. When the extraction theorem needs subexponentially many input copies, existing direct-sum copies are grouped before that step. This costs only $o(N)$ per round. Since there are finitely many rounds, the total overhead remains subexponential, and their interface tolerances can be chosen as positive constants independent of $N$.

\subsection{The revised optimization problem}

The rectangular asymptotic sum inequality uses the full volume of each output tensor \citep{schonhage1981}. Define
\begin{equation}\label{eq:mbar}
\Mbar=\frac13\sum_{W\in\WW}\sum_{T\in\mathcal L}M_{T,W}.
\end{equation}
This is the logarithm of the geometric mean of the three matrix sides. This volume-based evaluation follows directly from the asymptotic sum inequality.

\begin{resultbox}
\begin{theorem}[Final assembly]\label{thm:assembly}
For the level-four construction, with $\Mbar>0$ and the positive round rates required by the extraction theorem, every feasible solution of
\begin{equation}\label{eq:new-opt}
\begin{aligned}
\operatorname{minimize}\quad &\Omega\\[3pt]
\text{\normalfont subject to}\quad &
\underbrace{E_G+\Eshared-\frac DK}_{\text{\upshape copies}}
+\Omega\Mbar\ge8\log(q+2)
\end{aligned}
\end{equation}
provides the bound
\begin{equation}\label{eq:final-bound}
\omega\le\frac{8\log(q+2)-E_G-\Eshared+D/K}{\Mbar}.
\end{equation}
\end{theorem}
\end{resultbox}

\begin{proof}
For fixed $K$, the sixfold $K$-batch source has logarithmic border-rank budget at most $6KNC+o(N)$. Root extraction and \cref{eq:pipeline-total} produce at least
$2^{6KN(E_G+\Eshared-D/K-\rho)-o(N)}$ identical rectangular matrix multiplication summands. Their log volume is at least $3\cdot6KN(\Mbar-\rho)-o(N)$. The rectangular asymptotic sum inequality followed by $\rho\to0$ yields \cref{eq:final-bound}.
\end{proof}

The final bound combines the shared retained rate, rectangular output volume, and exact finite-batch accounting.

\section{Numerical optimization}\label{sec:optimization}

The shared extraction theorem rewards complementary imbalances across levels. We optimize its coupled objective to find parameters whose aggregate role capacities are nearly balanced.

The recorded search uses the programs and trajectory described in \cref{app:search}. It seeks a feasible point with enough margin for rational certification. The resulting rational witness supplies the parameters used in the theorem.

\subsection{Search objective and parameterization}

For fixed $q=5$ and $\ell^*=4$, let $\theta$ collect the continuous search variables. We first minimize the steady-state objective
\begin{equation}\label{eq:search-objective}
\Omega_\infty(\theta)=
\frac{C-E_G(\theta)-\Eshared(\theta)}{\Mbar(\theta)}.
\end{equation}
The finite construction is evaluated separately as
\begin{equation}\label{eq:finite-objective}
\Omega_K(\theta)=\Omega_\infty(\theta)+\frac{D(\theta)}{K\Mbar(\theta)}.
\end{equation}
The optimization trace records $\Omega_\infty$. Certification evaluates $\Omega_K$, including an outward bound on the finite-pipeline correction.

\paragraph{Probability coordinates.}
Simplex-valued parameters are represented by masked softmax coordinates. The parameter map omits inadmissible outcomes and fixes rows with only one admissible outcome. Complementary split entries are tied so that $\alpha_T^{(r)}(u)=\alpha_T^{(r)}(s_T-u)$, and terminal parameters are represented in the interval $(0,1/2)$. The compact parameter map is affine in the retained logit increments. It removes masked coordinates and ties the half-swap directions, while retaining the harmless additive softmax gauge directions for diagonal scaling.

\paragraph{Context and symmetry.}
The level-four source has $153$ shapes, of which $105$ are positive. Its positive children give $945$ positive level-three parent--child contexts. Each context carries six weighted strategies with its own split and child distributions. Distinct contexts are retained even when their shapes agree.

Complete-word distributions are represented on orbits of the prescribed binary-tree symmetries. At level four, this reduces $6{,}561$ words to $231$ orbits. If $p_O$ is an orbit's total mass and the distribution is uniform within the orbit, its entropy is evaluated as
\begin{equation}\label{eq:orbit-entropy}
H(\beta)=
\underbrace{H\bigl((p_O)_O\bigr)\vphantom{\sum_O p_O\log|O|}}_{\text{\upshape orbits}}
+\underbrace{\sum_O p_O\log|O|}_{\text{\upshape within orbits}}.
\end{equation}
The second term accounts for the uniform distribution within each orbit. The search is restricted to distributions with these prescribed symmetries.

\paragraph{Role allocation.}
The final recorded run starts from a prepared level-four point and keeps its selected level-four and level-three orientation allocations fixed. It optimizes the remaining continuous parameters, including terminal role allocations. For terminal nodes, write
\begin{equation}\label{eq:terminal-pooling}
B=\sum_Tm_T,\qquad
P=\sum_Tm_T(H_T-1)_+,\qquad
Q=\sum_Tm_T(1-H_T)_+.
\end{equation}
The implementation uses two probability vectors $v^+,v^-\in\Delta(\WW)$ and the terminal role vector
\begin{equation}\label{eq:terminal-two-policies}
e_2=B\mathbf1+P v^+-Q v^-.
\end{equation}
This corresponds to one explicit routing policy for nodes with $H_T\ge1$ and another for nodes with $H_T<1$. This is a restricted instance of the sector allocations in \cref{eq:terminal-role}.

\subsection{Conservative treatment of maximum entropy}\label{sec:gibbs}

The maximum-entropy terms in \cref{eq:hmax} are nested optimization problems. Because they enter the retained rate with a minus sign, a conservative evaluation requires upper bounds. Finite dual potentials provide these bounds at every iterate.

Let $\alpha$ be supported on $\mathcal D$, and choose strictly positive potentials $U_W$ for the three coordinates. Define
\begin{equation}\label{eq:gibbs}
T(u)=\prod_{W\in\WW}U_W(u_W),\qquad Z=\sum_{u\in\mathcal D}T(u).
\end{equation}
Then
\begin{equation}\label{eq:gibbs-bound}
H_{\mathcal D}^{\max}(\alpha)
\le \log Z-\sum_{u\in\mathcal D}\alpha(u)
       \sum_{W\in\WW}\log U_W(u_W).
\end{equation}
Indeed, every $\alpha'$ with the same marginals has the same expectation of $\log T(u)$, and the nonnegativity of its relative entropy with respect to $T/Z$ gives the inequality. Replacing $H^{\max}$ by the right-hand side can only decrease the first local retention rate. The potentials may be selected numerically and then rounded to positive rationals without changing the direction of the certificate.

\subsection{A constrained quasi-Newton search}

The objective is nonsmooth when two axes have equal retained rates. Near-equality of the final shared columns makes these points central to the search. The recorded solver models all competing axes explicitly.

For a trial exponent $\Omega_0$, write the steady-state feasibility residual as
\begin{equation}\label{eq:merit}
\Psi_{\Omega_0}(\theta)
=\sum_{j=1}^{J}\min_{W\in\WW}f_{j,W}(\theta)
 +\Omega_0\Mbar(\theta)-C.
\end{equation}
The groups $f_j$ contain the weighted root-region retention vectors and one shared constituent vector $e_4+e_3+e_2$. The recorded run uses $\Omega_0=2.371176$. When $\Mbar>0$, a nonnegative value of \cref{eq:merit} is equivalent to $\Omega_\infty\le\Omega_0$.

With a positive-definite local curvature model $B_\theta$, the sequential quadratic programming step has the form
\begin{equation}\label{eq:sqp}
\begin{aligned}
\underset{d,z}{\operatorname{maximize}}\quad&
  \sum_{j=1}^{J}z_j+
  \Omega_0\inner{\nabla\Mbar(\theta)}{d}
  -\tfrac12 d^\top B_\theta d,\\
\text{\normalfont subject to}\quad&
 z_j\le f_{j,W}(\theta)+\inner{\nabla f_{j,W}(\theta)}{d},\\
& j\in[J],\quad W\in\WW.
\end{aligned}
\end{equation}
Its dual has one three-entry probability vector per group. An inverse L-BFGS approximation with twelve curvature pairs supplies products with $B_\theta^{-1}$ in a memory-efficient representation.

The implementation uses JAX automatic differentiation and evaluates vector--Jacobian products sequentially to control peak memory. A diagonal, probability-dependent scaling compensates for the very different masses of the branches. Active axes are identified from the dual weights. The step is projected toward the tangent of their equality constraints, clipped in logit coordinates, and corrected for nonlinear equality error during backtracking. Acceptance requires an improvement in the actual residual \cref{eq:merit}.

These numerical techniques guide the search. The final proof depends on the feasibility of the selected rational witness.

\subsection{Recorded trajectory and stopping status}

Our optimization ends at
\begin{equation}\label{eq:search-endpoint}
\Omega_\infty=2.3710446603947700.
\end{equation}

The final accepted step improves the objective by approximately $1.03\cdot10^{-12}$. Global optimality remains unverified. The selected point has enough margin for rationalization and the finite-pipeline correction; exact certification details appear in \cref{app:formal}.

\section{Discussion}\label{sec:discussion}

We obtain the improvement by applying shared extraction to $\CW_5^{\otimes8}$, the source also used in the level-four construction of Dupont et al. Our local distributions, entropy penalties and terminal dimensions are inherited from the asymmetric laser framework. The mathematical change is to make factors at different levels share the collision budget of one extraction, and to realize that shared step in a finite schedule.

The new objective rewards a different balance of parameters. The extraction theorem makes complementary capacities usable in one round, and numerical optimization finds a witness that realizes the resulting improvement.

Larger strategy families, more flexible role allocations, higher CW powers, and pipelines containing more levels provide directions for further search. The attainable gains remain to be determined. The present bound concerns asymptotic complexity; the construction can require extremely large dimensions before its exponent becomes relevant.

\section{Acknowledgements}

This research is funded by the NSF under IIS Award 2552008. Computational resources for optimizing the bound and lean verification come from the \cite{Ohio_Supercomputer_Center1987-dl}.

\clearpage
\bibliographystyle{abbrvnat}
\bibliography{refs}

\clearpage
\appendix
\crefalias{section}{appendix}
\crefname{appendix}{appendix}{appendices}
\Crefname{appendix}{Appendix}{Appendices}
\section{Local formulas and notation}\label{app:local}

This appendix records the unchanged local ingredients in the notation used throughout the paper. All formulas use base-two entropy. For a nonnegative vector $v$, write
\begin{equation}\label{eq:perspective-entropy}
\PH(v)=\begin{cases}
\bigl(\sum_i v_i\bigr)H\bigl(v/\sum_i v_i\bigr),&\sum_i v_i>0,\\
0,&v=0.
\end{cases}
\end{equation}
This is simply the entropy of a pooled distribution multiplied by its mass, as in the pooled complete-split formulas of \citet{dupont2026}.

\subsection{Compatibility entropies at a positive node}

Fix a positive node $T$ of level at least three and the identity region. To shorten the formulas, let $s=s_T$, $\alpha=\alpha_T^{(1)}$, $b_{u,W}=\beta_{T[u,1],W}$, and
\[
w(u)=\alpha(u)+\alpha(s-u).
\]
The regional parent distribution is
$\beta_{T,W}^{(1)}=\sum_u\alpha(u)(b_{u,W}\times b_{s-u,W})$.
The compatibility entropies are
\begin{align}
\eta_{T,Y}^{(1)}
&=\sum_{u:u_Z=0}w(u)H(b_{u,Y})
 +\sum_j\PH\!\left(\sum_{\substack{u:u_Y=j\\u_Z>0}}w(u)b_{u,Y}\right),\label{eq:eta-Y}\\
\eta_{T,Z}^{(1)}
&=\sum_{u:u_X=0\,\text{or}\,u_Y=0}w(u)H(b_{u,Z})
 +\sum_k\PH\!\left(\sum_{\substack{u:u_Z=k\\u_X>0,\,u_Y>0}}w(u)b_{u,Z}\right).
\label{eq:eta-Z}
\end{align}
In each pooled term, all words have the same length and coordinate sum, so the vectors being added belong to the same space. The first term isolates the child sectors on which a zero coordinate forces a fine word from the earlier axis. The second term retains only the pooled law on the remaining sectors. The analogous formulas in region $r$ are obtained by applying $\pi_r$ to all axis labels.

\subsection{Root extraction}

Fix the identity region at the root. Put $\alpha=\alpha_G^{(1)}$ and $b_{s,W}=\beta_{G[s,1],W}$, for $s\in\Shapes_{\ell^*}$. Each root position contributes one child. Define
\begin{equation}
\bar\beta_{G,W}^{(1)}=\sum_s\alpha(s)b_{s,W}.
\end{equation}
The pooled root entropies are
\begin{align}
\eta_{G,Y}^{(1)}
&=\sum_{s:s_Z=0}\alpha(s)H(b_{s,Y})
 +\sum_j\PH\!\left(\sum_{\substack{s:s_Y=j\\s_Z>0}}\alpha(s)b_{s,Y}\right),\\
\eta_{G,Z}^{(1)}
&=\sum_{s:s_X=0\,\text{or}\,s_Y=0}\alpha(s)H(b_{s,Z})
 +\sum_k\PH\!\left(\sum_{\substack{s:s_Z=k\\s_X>0,\,s_Y>0}}\alpha(s)b_{s,Z}\right).
\end{align}
Thus
\begin{equation}\label{eq:root-rate}
E_G^{(1)}=\min\!\left\{
\begin{aligned}
&H(\alpha_X)-P_{\Shapes_{\ell^*}}(\alpha),\\
&H(\bar\beta_{G,Y}^{(1)})-\eta_{G,Y}^{(1)},\\
&H(\bar\beta_{G,Z}^{(1)})-\eta_{G,Z}^{(1)}
\end{aligned}\right\},
\qquad E_G=\sum_{r=1}^6A_G^{(r)}E_G^{(r)}.
\end{equation}
The other regions again follow by relabelling. The root contribution retains its regionwise minima.

\subsection{Leaf matrix sizes}

At a zero-shape leaf, let $W_0$ be its first zero coordinate and $W_1$ its first nonzero coordinate. The prescribed complete distribution gives
\begin{equation}\label{eq:zero-M}
M_{T,W_0}=m_T\!\left[
H(\beta_{T,W_1})+
\sum_L\beta_{T,W_1}(L)\,\#\{p:L_p=1\}\log q
\right],
\end{equation}
with the two other side logarithms equal to zero in this orientation.
For a positive terminal node with shape $(1,1,2)$,
\begin{equation}\label{eq:terminal-M}
(M_{T,X},M_{T,Y},M_{T,Z})
=m_T\log q\,(1-2\mu_T,\;1-2\mu_T,\;2\mu_T),
\end{equation}
and the coordinates are permuted for the other terminal shapes. These leaf dimensions enter the bottom-up and top-down evaluation unchanged.

\begin{table}[t]
\centering\small
\begin{tabularx}{\linewidth}{@{}lX@{}}
\toprule
Notation & Meaning\\
\midrule
$G$, $T$, $s_T$, $\ell$ & Root, node, node shape, and recursion level\\
$r$, $\pi_r$ & Region and its order of processing the physical axes\\
$A_T^{(r)}$, $\alpha_T^{(r)}$ & Region mass and distribution of left-child shapes\\
$m_T$ & Node multiplicity per unsymmetrized source copy\\
$\beta_{T,W}^{(r)}$, $\beta_{T,W}$ & Regional and mixed complete-word distributions\\
$P_D$, $\eta_{T,W}^{(r)}$ & Maximum-entropy penalty and pooled compatibility entropy\\
$E_{T,W}^{(r)}$, $E_\ell$, $E_G$ & Local retention, original level retention, and root retention\\
$e_{\ell,W}$, $\Eshared$ & Role-summed level capacities and their shared minimum\\
$M_{T,W}$, $\Mbar$ & Leaf log side and mean of the three total log sides\\
$K$, $D$, $\Omega$ & Fixed batch count, boundary loss, and candidate exponent\\
\bottomrule
\end{tabularx}
\caption{Notation follows the tree formulation of Dupont et al.; $e_{\ell,W}$, $\Eshared$ and $D$ describe the new assembly.}
\label{tab:notation}
\end{table}

\section{The heterogeneous extraction theorem}\label{app:theorem}

We give the precise statement used in the pipeline, keeping the underlying tensor and complete-type conditions explicit. Tensors are trilinear forms over the ground field. We write $T\dg T'$ when $T'$ is a degeneration of $T$; a restriction is the special case of constant linear maps. Direct sums use disjoint variables. Identifying corresponding variables maps a direct sum of tensors to their ordinary sum.

The Coppersmith--Winograd tensor is
\begin{equation}\label{eq:CW}
\begin{aligned}
\CW_q={}&\sum_{i=1}^q(x_0y_iz_i+x_iy_0z_i+x_iy_iz_0)\\
&+x_0y_0z_{q+1}+x_0y_{q+1}z_0+x_{q+1}y_0z_0.
\end{aligned}
\end{equation}
It has border rank at most $q+2$ \citep{cw1990}. Give variables indexed by $0$, by $1,\ldots,q$, and by $q+1$ degrees $0$, $1$ and $2$, respectively. Every nonzero term has total degree two. Consequently, at every base position in a tensor power, the three fine-word letters satisfy $x_p+y_p+z_p=2$.

A level-$\ell$ constituent $T_s^{(\ell)}$ is the part of $\CW_q^{\otimes 2^{\ell-1}}$ with shape $s\in\Shapes_\ell$. Its fine blocks are indexed by complete words, while its coarse blocks retain only their sums.

\subsection{Typed input and prescribed children}

Fix a finite set of types $\mathcal T$, independent of $N$. Type $t$ has level $\ell_t\ge2$, positive shape $s_t$, and positive rational mass $m_t$. Set
\[
n_t=m_tN,\qquad c_t=2^{\ell_t-1}.
\]
A parent of this type has $c_t$ base positions and two children with $c_t/2$ positions each. Choose a rational split $\alpha_t$ on $\Split(s_t)$ satisfying
\[
\alpha_t(u)=\alpha_t(s_t-u).
\]
For every positive-mass child and axis $W$, choose a rational complete distribution
$\beta_{t,u,W}\in\Delta(\mathcal C_{\ell_t-1,u_W})$. A zero coordinate has its all-zero word. If a child has a zero coordinate, the two remaining distributions must be related by the coordinatewise complement map $L\mapsto\vec2-L$.

Define
\begin{align}
\beta_{t,W}&=\sum_u\alpha_t(u)
   \bigl(\beta_{t,u,W}\times\beta_{t,s_t-u,W}\bigr),\label{eq:typed-beta}\\
\gamma_{t,W}&=(\alpha_t)_W,\qquad
P_t=H^{\max}_{\Split(s_t)}(\alpha_t)-H(\alpha_t),\qquad
w_t(u)=2\alpha_t(u).
\end{align}
The typed compatibility entropies $\eta_{t,Y},\eta_{t,Z}$ are \cref{eq:eta-Y,eq:eta-Z} with $w=w_t$ and $b_{u,W}=\beta_{t,u,W}$. Write
\begin{equation}\label{eq:typed-rates}
\begin{aligned}
e_{t,X}&=H(\gamma_{t,X})-P_t,\\
e_{t,Y}&=H(\beta_{t,Y})-\eta_{t,Y},\\
e_{t,Z}&=H(\beta_{t,Z})-\eta_{t,Z},
\qquad g=\min_{W\in\WW}\sum_t m_t e_{t,W}.
\end{aligned}
\end{equation}
A type in the tree formulation is a node--region pair, together with any strategy and sector labels that must be retained. Its weighted contribution is $m_te_{t,W}=E_{T,\pi_r(W)}^{(r)}$. Distinct contexts and sector labels are retained throughout the construction.

Let $P_N(\varepsilon)$ be
\[
\bigotimes_t\bigl(T_{s_t}^{(\ell_t)}\bigr)^{\otimes n_t},
\]
restricted on each axis to parent fine blocks whose empirical complete distributions are within $\varepsilon$ in sup norm of $\beta_{t,W}$. The exact child target is
\[
Q_N(0)=\bigotimes_{t,u}\bigl(T_u^{(\ell_t-1)}\bigr)^{\otimes n_tw_t(u)},
\]
restricted to the exact child complete types $\beta_{t,u,W}$. The approximate target $Q_N(\delta)$ allows sup-norm error $\delta$ in these types. Zero masses are omitted. We take $N$ through multiples for which every required exact population is integral.

\begin{theorem}[Heterogeneous extraction, precise form]\label{thm:heterogeneous-precise}
Assume $g>0$. For all sufficiently small fixed $\varepsilon>0$, $P_N(\varepsilon)$ degenerates into a direct sum of at least
\[
2^{N(g-r(\varepsilon))-o(N)}
\]
copies of $Q_N(0)$, where $r(\varepsilon)\to0$ as $\varepsilon\to0$.

For every $\rho>0$, there is $\varepsilon_0>0$ such that, for every $0<\varepsilon<\varepsilon_0$, there is $\delta_0>0$ with the following property. For $0<\delta<\delta_0$, subexponentially many independent copies of $P_N(\varepsilon)$ degenerate into at least $2^{N(g-\rho)-o(N)}$ copies of $Q_N(\delta)$. One may choose $\delta_0\le\varepsilon/4$.
\end{theorem}

The type list and the number of extraction rounds are fixed before $N$ grows. The pipeline in \cref{sec:shared} uses this quantifier order.

\section{Proof of heterogeneous extraction}\label{app:proof}

The proof extends the asymmetric extraction architecture of \citet{alman2025} by retaining the position type throughout counting, compatibility, and repair. Each step below establishes the required property for the heterogeneous product.

\subsection{Coarse types and the shared hash}

Encode a coarse triple by its left-child indices $(I_p,J_p,K_p)$, with $I_p+J_p+K_p=c_t$ at a position of type $t$. Keep vertices with the prescribed marginal types $\gamma_{t,W}$. Let $B_W$ count these vertices, $\Ngood$ the triples with joint types $\alpha_t$, and $\Nadm$ all triples with the retained marginals. Multinomial estimates give
\begin{align}
\log B_W&=N\sum_t m_tH(\gamma_{t,W})+O(\log N),\label{eq:count-vertices}\\
\log\Ngood&=N\sum_t m_tH(\alpha_t)+O(\log N),\label{eq:count-good}\\
\log\Nadm&\le N\sum_t m_tH_{\Split(s_t)}^{\max}(\alpha_t)+o(N).
\label{eq:count-all}
\end{align}
There are only polynomially many possible joint empirical types. The extraction uses the upper bound in \cref{eq:count-all}.

Permutations of parent positions within each type act transitively on each retained vertex class. Every $X$ vertex therefore has $\Nadm/B_X$ admissible incident triples, while every $Y$ or $Z$ vertex has $\Ngood/B_Y$ or $\Ngood/B_Z$ good incident triples. Counting in the unrestricted coarse support can only overestimate later collisions.

Use the hash of \cref{eq:hash} with an odd prime modulus $M$ larger than the fixed coarse alphabets. Keep hashes in a progression-free subset $\mathcal B\subset\ZZ/M\ZZ$ of size $M^{1-o(1)}$ \citep{behrend1946,alman2025}. Since $h_X+h_Y=2h_Z$, every surviving triple has a common bucket $b\in\mathcal B$.

For a fixed triple and bucket, every choice of the coefficients $w_p$ determines exactly one pair $b_0,w_0$ that puts all three hashes at $b$. Thus the probability is $M^{-2}$, and conditional on that event the $w_p$ remain independent and uniform. A distinct triple sharing the $X$ vertex collides when
\[
\sum_p w_p(J'_p-J_p)=0\pmod M.
\]
Its coefficient vector is nonzero modulo $M$, giving conditional probability $1/M$. The same argument applies to a shared $Y$ or $Z$ vertex. None of these identities requires a constant $c_t$.

Delete every $X$ vertex incident to more than one surviving triple, and every one whose unique triple has the wrong joint type. Conditional on a fixed good triple hashing to $b$, the probability of losing its $X$ vertex is at most $\Nadm/(B_XM)$. Every remaining $X$ vertex identifies one good triple, so its fine blocks may now be restricted to that triple's prescribed child types. This is an $X$-only variable deletion.

\subsection{Compatibility zero-outs}

For a good triple $e$, let $S_{t,u}(e)$ be its child positions of type $u$, taking both halves together. Its size is $n_tw_t(u)$. A fine block is \emph{useful} for $e$ when its complete distribution on each such set is $\beta_{t,u,W}$.

On $Y$, first impose the prescribed pooled word distribution on
$\bigcup_{u:u_Y=j}S_{t,u}(e)$ for each $(t,j)$. This set is determined by the coarse $Y$ vertex, hence by the fine $Y$ word, so that axis alone determines the restriction. Declare a fine $Y$ block compatible with $e$ if it satisfies these pooled laws and also has the individual law $\beta_{t,u,Y}$ whenever $u_Z=0$.

After $X$ usefulness, every surviving monomial's $Y$ block is compatible with its coarse triple. In a zero-$Z$ child, the $X$ and $Y$ words are complements; the imposed $X$ type therefore forces the required $Y$ type. This implication holds locally for each child length.

Delete a fine $Y$ block when it is compatible with several remaining good triples. A survivor now identifies its unique triple and can be restricted to full usefulness. For $Z$, proceed in the same order: pooled laws for each $(t,k)$, followed by individual laws in sectors with $u_X=0$ or $u_Y=0$, forced respectively by $Y$ or $X$ usefulness; then multiple-compatibility deletion and full usefulness.

The remaining triples have disjoint $X$ variables and disjoint fine $Y$ and $Z$ blocks. Their tensor is a direct sum of restrictions of the intended child products. The order of these operations is essential: restrictions requiring the full coarse triple are imposed only after the relevant axis identifies it uniquely.

\subsection{Double counting compatible blocks}

Fix a fine $Y$ block with parent empirical types $\xi_{t,Y}$ at coarse vertex $J$. Let $p_Y(\xi)$ be the fraction of good triples incident to $J$ with which the block is compatible. Permutations within each type and coarse-$Y$ pair class act transitively on blocks with these empirical types, so this fraction depends only on $\xi$.

Count pairs $(e,f)$ where $e$ is good and $f$ is a fine $Y$ block in its vertex with types $\xi$. Their count has logarithm
\begin{equation}\label{eq:pair-count}
\log\Ngood+\sum_t n_t\bigl[H(\xi_{t,Y})-H(\gamma_{t,Y})\bigr]+O(\log N).
\end{equation}
For a fixed $e$, discard the parent empirical-type constraint to upper-bound the compatible $f$. The zero-$Z$ child sectors have individual prescribed laws. On the other sectors with the same $Y$ index, subtracting the isolated sectors from the prescribed total leaves a pooled law. Their combined multinomial exponent is exactly $\sum_t n_t\eta_{t,Y}$. Dividing the compatible-pair upper bound by \cref{eq:pair-count} gives
\begin{equation}\label{eq:compatibility-bound}
p_Y(\xi)\le
2^{\sum_t n_t[\eta_{t,Y}-H(\xi_{t,Y})+H(\gamma_{t,Y})]+O(\log N)}.
\end{equation}
Finite-alphabet entropy continuity replaces $\xi_{t,Y}$ by $\beta_{t,Y}$, uniformly over the parent tolerance window, at cost $Nr_Y(\varepsilon)$ with $r_Y(\varepsilon)\to0$. The same count yields
\begin{equation}
p_Z\le
2^{\sum_t n_t[\eta_{t,Z}-H(\beta_{t,Z})+H(\gamma_{t,Z})]
+Nr_Z(\varepsilon)+O(\log N)}.
\end{equation}
Both bounds can be capped at one. Each calculation concerns the marginal distribution on one axis and permits dependencies between the three axis words.

A fixed typical fine $Y$ block has at most $(\Ngood/B_Y)p_Y$ compatible good competitors sharing its coarse vertex. Conditional on its own triple hashing to $b$, each competitor collides with probability $1/M$. The union bound therefore gives deletion probability at most $\Ngood p_Y/(B_YM)$, and analogously $\Ngood p_Z/(B_ZM)$ for $Z$. Counting all potential competitors makes the union bound valid with dependencies between stages.

\subsection{Concentration and the number of survivors}

Fix a good coarse triple and its exact child target. On one axis, a uniformly random target fine block assigns the prescribed multiset of complete words uniformly within each child-type group. For two different child groups, a parent pair has the product word law in \cref{eq:typed-beta}. If both child types agree, the two words are sampled without replacement from the same group. For a group of $r$ slots, the pair probability for words $a,b$ is
\[
\frac{[r\beta(a)]\,[r\beta(b)-\mathbf1_{a=b}]}{r(r-1)},
\]
which differs from $\beta(a)\beta(b)$ by $O(1/r)$. All positive-mass groups have $r=\Theta(N)$. The expected parent empirical distribution is therefore within $O(1/N)$ of the prescribed one.

Exposing the multiset permutations gives a bounded-differences martingale: a transposition changes at most two parent words. Consequently, for fixed $\varepsilon>0$, the fraction of target blocks outside the parent tolerance window is at most $C_1e^{-c\varepsilon^2N}$. The estimate applies separately to the uniform fine-block distribution on each axis. The same argument is uniform for nearby child laws.

Choose $M$ using \cref{eq:modulus}, with its polynomial slack $N^8$. Each conditional collision-deletion probability is then at most $N^{-8}$. Linearity of expectation over fine blocks and Markov's inequality show that, conditional on a good triple hashing to a bucket, it survives with $O(N^{-2})$ holes on every axis with probability $1-O(N^{-6})$. Averaging over good triples and buckets supplies a hash choice with at least
\begin{equation}\label{eq:broken-copies}
(1-o(1))\Ngood\,|\mathcal B|/M^2=\Ngood M^{-1-o(1)}
\end{equation}
independent broken copies, all intended to be the same child tensor.

Substituting the counting and compatibility bounds into the choice of $M$ gives
\begin{equation}\label{eq:proof-rate}
\log\Ngood-\log M\ge
N\min_{W\in\WW}\sum_t m_te_{t,W}
-Nr(\varepsilon)-o(N).
\end{equation}
Since $\Nadm/B_X\ge1$, the maximum with one in the modulus has the same exponential rate. Equation~\eqref{eq:proof-rate} is the desired shared rate, before repairing the holes.

\subsection{Repair for a heterogeneous target}

The repair step uses the partitioned-tensor principle of \citet[Section~7]{vxxz2024}. We verify its required action for the present mixed-level product.

Suppose a tensor is partitioned into $n_X,n_Y,n_Z$ parts and has an automorphism group preserving the partitions and transitive on the parts of each axis. Each available broken copy is missing at most a fraction $h$ of the parts on every axis, with arbitrary hole patterns. For a subproblem on part sets $A_X,A_Y,A_Z$, a random group element moves the holes into each $A_W$ with expected intersection at most $h|A_W|$. Markov's inequality and a union bound give a group element for which all three intersections have size at most $4h|A_W|$.

A fresh moved copy supplies the box of nonmissing parts. The remainder consists of seven axis boxes, each using a missing subset on at least one axis. Each child's product of part counts is at most $4h$ times its parent's. Recursion with fresh copies has depth at most
\begin{equation}
1+\left\lceil\frac{\log(n_Xn_Yn_Z)}{\log(1/(4h))}\right\rceil.
\end{equation}
The resulting boxes partition the target monomials, so identifying corresponding variables supplies each target coefficient exactly once.

For $Q_N(0)$, independently permute child-factor positions within each fixed $(t,u)$. These permutations preserve the tensor and exact-type conditions, map whole fine blocks to whole fine blocks, and act transitively on each axis's fine blocks. Separate transitivity on each axis supplies the required marginal uniformity under a single group action. Since there are $\exp(O(N))$ fine blocks and $h=O(N^{-2})$, the depth is $O(N/\log N)$. The seven-way recursion therefore requires only $2^{O(N/\log N)}=2^{o(N)}$ broken copies per repaired exact copy. Grouping the copies in \cref{eq:broken-copies} preserves the rate in \cref{eq:proof-rate}.

\subsection{Exact to approximate output}

The next extraction uses an approximate interface. For a fixed finite list of alphabets, $Q_N(\delta)$ is the ordinary sum of at most $(N+1)^{D_0}$ exact-type subtensors, for a constant $D_0$. Every monomial belongs to a unique tuple of empirical types. Tuples incompatible with tensor support or zero-coordinate complementarity contribute zero and can be omitted.

For any admissible child tuple $\beta'$ within $\delta$ of $\beta$, form its parent laws using \cref{eq:typed-beta}. The elementary bound $|ab-a'b'|\le|a-a'|+|b-b'|$ for probability entries gives
\[
\norm{\beta'_{t,W}-\beta_{t,W}}_\infty\le2\delta.
\]
Thus, when $\delta\le\varepsilon/4$, the parent interface centered at $\beta'_{t,W}$ with tolerance $\varepsilon/2$ is a restriction of $P_N(\varepsilon)$.

Apply the exact-output result to one independent input copy per admissible exact tuple, retain the common minimum number of outputs, and identify corresponding variables to sum the exact subtensors at each output index. Their monomials partition $Q_N(\delta)$, so the sum is precisely the desired approximate output. The additional input multiplicity is polynomial, hence subexponential. Entropy continuity and the uniform concentration estimate allow the losses to be split between a sufficiently small $\varepsilon$ and then a sufficiently small $\delta$. This proves the quantifiers in \cref{thm:heterogeneous-precise}. \qed

\section{Sector refinement and role routing}\label{app:roles}

This appendix realizes the optimizer's role allocations as tensor restrictions.

\subsection{Refining one strategy into sectors}

Suppose a parent complete distribution is a mixture
$\beta_W=\sum_s a_s\beta_W^s$. Divide the parent positions into fixed sectors with relative sizes $a_s$ and impose the distribution $\beta_W^s$ on sector $s$. The aggregate empirical distribution is the required mixture, so this is a restriction of the original interface, with the corresponding tolerance allowance. The sectors may remain separately labelled throughout subsequent extraction.

This realizes the six weighted strategies at each level-three context. Each strategy can then be allocated among the six asymmetric orders while retaining the same split and child distributions in every allocated sector. The strategy retains its prescribed descendants under every order, with the corresponding three role capacities. These are the allocations denoted by \code{alloc4} and \code{alloc3} in the certificate.

\subsection{Aligning physical orders}

Take six copies of the entire source construction, indexed by coordinate permutations $\sigma\in S_3$. Extract their roots separately. A strategy with order $\pi$ inside copy $\sigma$ has physical order $\sigma\pi$ and retains its capacity vector in role coordinates.

For each fixed physical order $\varphi$ and each original strategy order $\pi$, exactly one copy, $\sigma=\varphi\pi^{-1}$, places that strategy in physical order $\varphi$. Regrouping by physical order therefore creates six regions with identical sums of role capacities. Apply \cref{thm:heterogeneous-precise} within each region. The factor of six in the number of regions cancels against the six source copies, leaving the minimum of the role sum per original source.

The output is still the union of all coordinate-permuted, sector-labelled child lists. Its symmetry under coordinate permutations permits the same routing at the next round. Contexts and sectors remain separately labelled when their numerical contributions are summed. The root rate enters the final accounting separately.

\subsection{The terminal policy}

For shape $(1,1,2)$, the symmetric split has masses
\[
\mu,\quad\mu,\quad(1-2\mu)/2,\quad(1-2\mu)/2
\]
on children $(0,0,2),(1,1,0),(0,1,1),(1,0,1)$. Its marginals determine the joint split, so its entropy penalty is zero. The child complete words have length one, and the local role capacities are a permutation of
$(1,1,H(\mu,\mu,1-2\mu))$.

Assigning the doubled coordinate to role $W$ in proportion $v_W$ gives
\[
m\bigl[\mathbf1+(H(\mu,\mu,1-2\mu)-1)v\bigr].
\]
The policy is implemented by the same fixed-sector restriction and sixfold routing. This proves the interpretation of \cref{eq:terminal-role,eq:terminal-two-policies}.

\section{Certificate and formal verification details}\label{app:formal}

This appendix records the scope and implementation details of the accompanying formalization, as distinct from the mathematical exposition and the floating-point search.

\subsection{Exact parameter data}

All probability rows in the supplied certificate use denominator $2^{42}$. The root distribution ranges over $153$ level-four shapes. Each of the $105$ positive shapes has a split over the $45$ level-three shapes and an allocation among six role orders. There are $945$ positive parent--child contexts at level three; each has six weighted strategies, each with a split over the $15$ level-two shapes and an allocation among six orders. The terminal parameters have shape $945\times3\times6$.

\begin{table}[t]
\centering\small
\begin{tabularx}{\linewidth}{@{}l l X@{}}
\toprule
Data & Size or count & Role\\
\midrule
Root shapes & $153$ & Distribution of source constituents\\
\code{alpha4}, \code{alloc4} & $105\times45$, $105\times6$ & Positive level-four splits and routing\\
\code{alpha3} & $945\times6\times15$ & Contextual level-three strategy splits\\
\code{alloc3} & $945\times6\times6$ & Strategy routing among physical orders\\
Terminal parameters & $17{,}010$ & Level-two values of $\mu$\\
Simplex rows & $15{,}600$ & Exact probability normalization\\
Split rows & $5{,}770$ & Admissible, half-symmetric splits\\
Gibbs rows & $17{,}298$ & Positive maximum-entropy potentials\\
Zero-leaf rows & $8{,}879$ & Complete distributions for matrix leaves\\
\bottomrule
\end{tabularx}
\caption{Structure of the supplied exact certificate. The table counts the stored certificate objects.}
\label{tab:certificate-data}
\end{table}

Each simplex row sums exactly to its denominator; split rows have admissible support and complementary symmetry; every terminal $\mu$ lies in $(0,1/2)$; and all Gibbs entries are positive. The fifteen scalar rate expressions comprise three axes at each of the root, level four, level three, terminal conversion and matrix dimension. They expand into $234{,}991$ weighted logarithms of $167{,}384$ distinct positive rationals, organized into $467$ blocks of at most $512$ terms.

The formal integer logarithm checker uses range reduction over $256$ proved grid values, a cubic approximation with a proved error bound, and rational outward endpoints at scale $2^{60}$. Interval evaluation of \cref{eq:boundary,eq:final-bound} then establishes positivity of the residual at $2371045/10^6$ and $K=10^6$. All numerical proof steps use exact rational arithmetic.

\subsection{Statement and scope}

The stated endpoint in the Lean source is
\begin{algorithmnote}
\small\ttfamily
\noindent theorem MatrixBounds.Numeric.SuppliedCertifiedPipeline.exponent\_lt\\
\hspace*{1em}\{K : Type\}\ [CommRing K] :\\
\hspace*{2em}MatrixComplexity.exponent K <\\
\hspace*{3em}(23710449 : $\mathbb R$) / 10000000
\end{algorithmnote}
The formal bound $2.3710449$ is slightly stronger than the rounded value $2.371045$ stated in \cref{thm:main}. Here \code{K} denotes the ground ring; the batch count uses the symbol $K$ elsewhere in the paper. Its \code{exponent} definition is exactly the rank-based quantity in \cref{eq:omega}. The theorem is stated for rings in \code{Type}. The interpretation of $\omega$ over fields additionally uses the conventional rank-to-operation-count equivalence, whose formalization lies outside this formalization.

The mathematical presentation in the main text uses a characteristic-zero field. The formal theorem applies to the exact-rank definition over the commutative rings specified above.

\subsection{Scope of the formal verification}

The formalization includes finite tensors, exact rank decompositions, direct sums and tensor products; polynomial degenerations and coefficient extraction; the $(q+2)$-term CW degeneration; and the asymptotic sum inequality for identical rectangular summands. For $q=5$ and level four, the source certificate has rank bound $7^8$ and degeneration degree $24$.

The extraction development checks the shared hash identity and conditional collision probabilities, progression-free buckets from Mathlib's Behrend construction, the necessity of compatibility tests for nonzero CW terms, concentration, repair, and approximate interfaces. Its concentration proof uses a second-moment estimate sufficient for subexponential repair; the exposition in \cref{app:proof} uses a stronger bounded-differences estimate for readability. The formal contextual statement preserves waiting tensor factors and earlier copy labels, which supports composition throughout the finite pipeline.

Sector allocation, sixfold routing, terminal conversion, and the $K+2$-round schedule then provide genuine rectangular algorithms with the asserted source-rank, copy-count and volume bounds for every positive asymptotic error allowance.

\paragraph{Binding the mathematical rates to the data.}
The construction is tied to the certificate by twelve axis-rate identities and one matrix-volume identity. The proof proceeds through root-coarse, root-fine, paired-coarse, paired-fine, terminal and dimension streams. On the construction side, the rates are sums of parent entropies minus isolated-child and pooled-sector entropies, or coarse Gibbs terms, weighted by integer populations. On the certificate side, they are fixed weighted-log term lists.

These expressions are normalized node by node and compared on exact integer numerators in source-aligned windows. Tabulated corrections for terms spanning neighboring windows are proved to cancel. The endpoint theorem applies the numerical inequality only after these identities and the actual extractions are established. This binds the checked scalar inequality to the tensor construction.

\subsection{Audit and independent reproduction}

The development uses Lean~4.24.0 and Mathlib commit \code{f897ebcf} \citep{lean4,mathlib2020}. It comprises $5{,}492$ root modules with $47{,}316$ named theorems and $88{,}134$ named declarations. The source contains no \code{sorry}, \code{admit}, \code{native\_decide} or custom \code{axiom}, and no user-defined syntax, macros or elaborators. Numeric comparisons use kernel evaluation on exact rationals. A set of generated audit modules queries the axioms of every named declaration.

\paragraph{Build.} A from-scratch build is $8{,}899$ Lake jobs. Many generated modules need $4$--$10$~GiB of memory each; on one $96$-core node with \code{LEAN\_NUM\_THREADS=48} the build takes about $90$ minutes. With the pinned repository, the build commands are
\begin{algorithmnote}
\small\ttfamily
\noindent lake exe cache get\\
lake build
\end{algorithmnote}

\paragraph{Independent reproduction.} The first author developed and built the formalization. The second author then rebuilt the released Lean sources (tree \code{bb7aadd4}, at the time commit \code{bd86c75}) from a fresh checkout, with a freshly installed toolchain and Mathlib cache, on a separate compute allocation: all $8{,}899$ jobs succeeded with no errors and no \code{sorry} warnings. Across the $88{,}168$ audit lines of that build, the only axioms that occur are \code{propext}, \code{Classical.choice} and \code{Quot.sound}, and these three are the dependencies reported for \code{exponent\_lt}. The \code{lean4checker} tool (tag \code{v4.24.0}), which replays every declaration of a compiled module through the Lean kernel to detect tampering with the environment, was then run on each of the $6{,}534$ source modules of that build, including the generated audit parts, and reported no errors. The standalone Python interval verifier was also rerun on the shipped certificate and encloses the finite-pipeline bound in $[2.37104473,\,2.37104474]$.

\paragraph{Statement audit.} The statement of \code{exponent\_lt} depends on a small trusted surface: the definitions of a rank decomposition, the matrix multiplication tensor, and \code{MatrixComplexity.exponent}, about $120$ lines that import only Mathlib. The development proves that the tensor's bilinear evaluation equals Mathlib's matrix product, that the set defining the exponent is nonempty, and that $0\le\omega_R\le3$. We provide a standalone challenge file that restates these definitions verbatim and the theorem with its proof omitted, so that the correspondence between the formal statement and \cref{eq:omega} can be checked without reading the proof. The second author audited this statement. We then ran the Lean FRO's Comparator tool against the challenge file: it builds the challenge and the project in separate sandboxes, exports the statement's dependency closure from each, checks that every constant on which the statement of \code{exponent\_lt} depends is identical in the two environments, checks that the exported proof uses only the three permitted axioms, and replays the entire exported proof (19.3 million export lines) through a second Lean kernel. All steps passed; the replay took $20.6$ hours single-threaded and $16.7$~GB of memory. Because no Comparator release targets Lean~4.24, we used the October 2025 commits of Comparator and \code{lean4export} that match that toolchain, together with \code{landrun} built from source; the scripts, configuration, challenge file and logs are archived with the development. Kernel checking establishes the theorem relative to these definitions and Lean's trusted foundation; we encourage further independent review of the definitions.

\section{Search artifacts and reproducibility notes}\label{app:search}

\subsection{Search procedure}

The numerical search used in the final stage is summarized in \cref{alg:search}. It operates entirely as a discovery procedure; the theorem depends only on the certified rational witness produced afterward.

\begin{algorithm}[H]
\caption{Search for a shared-extraction witness}\label{alg:search}
\begin{algorithmic}[1]
\Require A feasible tree parameterization, fixed orientation allocations, trial $\Omega_0$
\State Encode the free simplex entries and tied split coordinates in logits.
\State Compute word distributions bottom-up and node masses top-down.
\While{continuing the numerical search}
  \State Evaluate $f_j$, $\Mbar$, their derivatives, and the residual $\Psi_{\Omega_0}$.
  \State Solve the small dual of \cref{eq:sqp} using inverse L-BFGS products.
  \State Project, clip, and backtrack the step; correct active-axis equalities.
  \State Accept only an improving step and retain the best evaluated witness.
\EndWhile
\State Round the selected witness to exact rationals and check its support and sums.
\State Restore $D/K$ and certify the final inequality independently of the solver.
\end{algorithmic}
\end{algorithm}

\subsection{Rate structure of the selected witness}

\Cref{tab:rates} shows the selected point's rate structure. Level three supplies relatively more capacity in role $X$ than in role $Z$; level two supplies the opposite pattern. Their complementary rows give almost equal shared sums. This is the complementarity that the new extraction makes usable.

\begin{table}[H]
\centering\small
\begin{tabular}{@{}lrrr@{}}
\toprule
Quantity & Role $X$ & Role $Y$ & Role $Z$\\
\midrule
$e_{4,W}$ & $1.663716411$ & $1.645107831$ & $1.631344714$\\
$e_{3,W}$ & $2.267823262$ & $2.208135648$ & $1.890628640$\\
$e_{2,W}$ & $1.509416740$ & $1.587712934$ & $1.918983059$\\
\rowcolor{wash}
$\sum_{\ell=2}^{4}e_{\ell,W}$ & $\mathbf{5.440956413}$ & $\mathbf{5.440956413}$ & $\mathbf{5.440956413}$\\
\midrule
\multicolumn{4}{@{}l@{}}{\textit{Other rates, in bits per source copy before symmetrization}}\\
$C=8\log7$ & \multicolumn{3}{r@{}}{$22.458839$}\\
$E_G$ & \multicolumn{3}{r@{}}{$2.656688$}\\
$\sum_\ell\min_W e_{\ell,W}$ & \multicolumn{3}{r@{}}{$5.031390095$}\\
$D$ & \multicolumn{3}{r@{}}{$[0.4419380147,\ 0.4419380157]$}\\
\midrule
\multicolumn{4}{@{}l@{}}{\textit{Matrix log sides, in physical coordinates}}\\
$\sum_{T\in\mathcal L}M_{T,W}$ & $6.097497$ & $6.021780$ & $6.051441$\\
$\Mbar$ & \multicolumn{3}{r@{}}{$6.056906$}\\
\bottomrule
\end{tabular}
\caption{The selected certificate's rates. Displayed decimals summarize the exact certificate data. The matrices' physical axes in the final two rows are distinct from the hashing roles in the first four.}
\label{tab:rates}
\end{table}

As a diagnostic, keeping this same point but replacing the shared retention by the sum of its levelwise minima gives approximately $2.43866$, even when retaining the mean matrix-side logarithm. This fixed-parameter diagnostic shows how strongly the selected point relies on shared extraction. We did not run a matched ablation. Two data points bound the comparison. Our own searches in the level-by-level formulation, before the shared extraction was introduced, reached approximately $2.37122$ and did not improve on \citet{dupont2026}. The transcript of \citet{naslund2026} reports $2.371161$ in that formulation. Both suggest that most of the improvement to $2.371045$ comes from the shared extraction and not from the search.

The optimization section uses three stored programs: \code{pipeline\_all\_sqp.py}, \code{pipeline\_context.py} and \code{compact\_context.py}. The trajectory is recorded in \code{pipeline\_all\_log.jsonl}. The source package includes these records and a CSV version of the trace. Reproducing the full optimization additionally requires its dependency modules and saved states; the Lean development is distributed as a separate artifact.

For precision, the following implementation details matter when interpreting the run. The inverse L-BFGS memory is twelve pairs. Dual weights above $10^{-6}$ identify active axes. The code alternates coordinate clipping at magnitude $1.5$ with tangent reprojection, then performs backtracking with up to four nonlinear equality corrections per trial. A candidate step must improve the actual merit by more than both $10^{-13}$ and $0.02$ times the predicted gain. Failed steps reset the curvature history and reduce the scale. These safeguards govern the numerical search. Exact certification supplies the final feasibility guarantee.

The search trajectory, outward interval certificate and Lean result are different artifacts. The first records a floating-point discovery process; the second certifies one scalar inequality under the extraction formulas; the third is the end-to-end exact-rank formalization. Each has its own timing, dependencies, and verification scope.

\section{Timeline}\label{sec:timeline}

\begin{description}
\item[September 17, $\sim$3:00 p.m.] Prompting began with OpenAI GPT-6 Astra. The initial prompt asked, in simple terms, for significant analytical progress toward improving the matrix-multiplication exponent bound.

\item[September 17, $\sim$7:00 p.m.] After approximately four hours, Astra produced a natural-language proof judged suitable for formalization.

\item[September 17--19] The Astra conversation was passed to Codex Astra xhigh with the instruction ``Formalize the proof.'' Codex worked on the Lean formalization for approximately 30 hours (mostly uninterrupted besides token limit resets) before available usage was exhausted.

\item[September 19] The incomplete Lean development was transferred to Claude Fable 5.1 in Ultra Code mode. After approximately eight additional hours, Fable completed the Lean certificate.

\item[September 19] Fable drafted the initial manuscript, which was subsequently edited with Astra.

\item[September 19] The first completed certificate established $\omega<2.371087$. Review of the numerical construction revealed that the optimization had not yet fully converged.

\item[September 19] Meta Muse Spark 1.3 was used on the Ohio Supercomputer Center to continue the optimization with GPU acceleration and tests of multiple initializations.

\item[September 19] The optimization converged to the final reported parameter set, improving the bound to $\omega<2.3710449$. The Lean certificate was regenerated using the updated parameters, and the manuscript was revised accordingly.

\item[September 19--21] With the draft complete, the authors began discussing the paper's layout and contents, corrections for accuracy, and additions needed for clarity and transparency.

\item[September 20--21] The second author, assisted by Claude Fable~5.1 in Claude Code, audited the formal statement, rebuilt the development from a fresh checkout on a separate allocation, ran \code{lean4checker} on every module, and revised the manuscript (\cref{app:formal}).

\end{description}

\paragraph{Statement on AI use and author responsibility.}
The timeline above records which system did what and when. In summary, OpenAI's GPT-6 Astra produced the shared-extraction idea, its natural-language proof (\cref{app:theorem,app:proof,app:roles}) and the first numerical search. Astra, running in Codex, and Anthropic's Claude Fable~5.1 wrote the Lean formalization. Meta's Muse Spark~1.3 ran the final optimization to convergence. Fable drafted the manuscript, which was then edited with Astra, with Fable and by the authors; Fable also assisted the second author with the independent verification. The authors chose the problem, directed the process, audited the formal statement, and independently reproduced the build (\cref{app:formal}). As stated in \cref{sec:preface}, we have not independently verified every step of the natural-language proof in \cref{app:proof}; for correctness we rely on the formal proof. We take responsibility for the theorem as formally stated and for the accuracy of this description.

\end{document}
